% !TEX program = lualatex
\documentclass[aspectratio=43,professionalfonts,handout]{beamer}
\usepackage{beamer_template}

\newcommand{\LectureNum}{1}
\newcommand{\LectureTitle}{ラプラス変換の定義}
\newcommand{\TextPages}{p.\,1-8}
\newcommand{\NextTopic}{相似性と移動法則}
\newcommand{\NextPages}{p.\,8-14}
\newcommand{\CourseName}{応用数学}
\renewcommand{\TermName}{前期}

\begin{document}

% -----------------------------------------------
% スライド1: 表紙＋目標＋用語・定義
% -----------------------------------------------
\begin{frame}{\CourseName\quad 第\LectureNum 回「\LectureTitle 」}
  {\small \TermName\quad \TeacherName\quad ／\quad 教科書 \TextPages\quad
  \textbf{目標}：ラプラス変換の定義を理解し、基本的な関数のラプラス変換を求めることができる}

  \begin{block}{用語・定義}
  \begin{description}[labelwidth=5.5em]
    \item[\kw{ラプラス変換}] ~
    \item[\kw{原関数 f(t)}] ~
    \item[\kw{像関数 F(s)}] ~
    \item[\kw{線形性}] ~
  \end{description}
  \end{block}
\end{frame}

% -----------------------------------------------
% スライド2: 公式・性質
% -----------------------------------------------
\begin{frame}{公式・性質}
  \begin{block}{}
  \begin{itemize}
    \item ラプラス変換の定義：$\mathcal{L}[f(t)] = F(s) = \int_0^\infty f(t)e^{-st}\,dt$
    \item $\mathcal{L}[1] = \frac{1}{s}$
    \item $\mathcal{L}[t^n] = \frac{n!}{s^{n+1}}$
    \item $\mathcal{L}[e^{at}] = \frac{1}{s-a}$
    \item 線形性：$\mathcal{L}[af(t)+bg(t)] = aF(s)+bG(s)$
  \end{itemize}
  \end{block}
\end{frame}

% -----------------------------------------------
% スライド3: 図・グラフ（TikZコードを手動追加）
% -----------------------------------------------
\begin{frame}{図・グラフ}
  \centering
  \vfill
  {\large （授業中に板書・図示）}
  \vfill
\end{frame}

% -----------------------------------------------
% スライド4: 例題→問
% -----------------------------------------------
\begin{frame}{例題→問}
  \begin{exampleblock}{例1) p.3}
    （板書で解説）
  \end{exampleblock}

  \vfill

  \begin{block}{問1) p.4　自分で解いてみよう}
    （ノートに解くこと）
  \end{block}
\end{frame}

% -----------------------------------------------
% スライド5: 例題→問
% -----------------------------------------------
\begin{frame}{例題→問}
  \begin{exampleblock}{例2) p.5}
    （板書で解説）
  \end{exampleblock}

  \vfill

  \begin{block}{問2) p.6　自分で解いてみよう}
    （ノートに解くこと）
  \end{block}
\end{frame}

% -----------------------------------------------
% まとめ
% -----------------------------------------------
\begin{frame}{まとめ}
  \begin{itemize}
    \item ラプラス変換の定義 $F(s)=\int_0^\infty f(t)e^{-st}\,dt$
    \item 基本関数のラプラス変換（$1, t, t^n, e^{at}$）
    \item ラプラス変換の線形性
  \end{itemize}
\end{frame}

\end{document}